\section{Philosophical Significance of Golden Reasoning}

\subsection{Agency's Curse}

To begin this paper, we would like to elaborate on an epistemological problem which operates in a similar vein to skepticism, however with different enough conditions so to as to warrant a separate definition. It will furthermore be argued that in lieu of the demotion with Aethic reasoning \cite{Aethus} of several key metaphysical problems to irrelevant at best, and attributions of constructs at worst, it now follows that the central problem to be addressed in this paper is perhaps humanity's highest remaining\footnote{Or, perhaps the highest \textit{now illuminated} issue.} philosophical issue. That is, we can argue this philosophical demon, which we shall call \textit{agency's curse}, is to the field of philosophy as philosophical skepticism is to the field of science. While skepticism was never truly dissolved, (nor could it ever be), we were able to substantially defang it a few hundred years ago with the development of the scientific method. In doing so, we created a clear distinction between that which is scientific and that which is ascientific \cite{}, and as such triggered massive leaps in our technological and metaphysical development as humans. However, as is somewhat striking, we are to argue that there is a close analogue to this duet of skepticism and the scientific method, but this time present in philosophy. Such a thing would then place us in the analogous history of philosophy as close to where Francis Bacon or René Descartes were in the history of science. At our location in history, then, an even greater epistemological demon\footnote{Of course not a demon in the literal sense -- instead in the metaphorical sense.} than skepticism now looms over humanity through the avenue of our philosophy. In the following section, let us go over the concept of \textit{agency's curse}, and after that argue how we might defang it. In doing so, we may be able to do for philosophy what the advent of science did to the middle ages almost five hundred years ago. 

To begin with, let us give a title to this central epistemological problem. Suppose we call it, \textit{agency's curse}. The idea, simply put, is this -- we, as agents in the universe, have two fundamentally conflicting properties to us.
\begin{enumerate}[label=Condition \arabic*.]
    \item We, at every moment, are bound to make actions of influence upon the universe. Even inaction, sleeping, or suicide are themselves all representations of actions which have effects on the universe around us as opposed to the alternative. In essence, then, we are in a continuous and vastly propagating flux of decisions at all moments through which the future of the entire universe depends. This is the metaphysical equivalent to a mother having to navigate an obstacle course of flames to save her child from a burning building, but the child keeps sliding ten more feet away. The stakes are infinitely broader than the time allotted for a decision, and the nature of the situation endlessly binds us to these decisions inherently. If we had any semblance of an idea about how many future eons of consequence we are either sloppily enabling or sloppily desecrating on a second to second basis, it would drive us to some kind of hypnotically overstimulated madness. By simply existing, we are permanently unable to escape this continuous and overbearing strain that is our own forced agency coupled with our own forced ignorance of the implications of our actions. It is Sisyphus, but with the evolving fate of the universe in place of a boulder. 
    \item We have no means of accessing the objective morality/meaning/interpretation of the universe, if there even is any. In a sense, then, if we picture a spectrum from nonliving matter to a deity, we are just high enough on this spectrum to possess agency and thus comply with the aforementioned condition, but we are just low enough to not be able to weigh our actions against any backdrop of metaphysically invariant reasoning. This is a very specific notion that we be elaborated upon in detail momentarily. Also, note that it is logical side effect we hold through the very nature of existing in the first place, rather than being some external property which is tagged onto humans in particular. 
\end{enumerate}
This second condition of agency can be highlighted more concretely in a coming thought experiment. But before we perform such a thought experiment, consider some vocabulary we will use to describe the second condition of agency's curse. Firstly, there will be a ``\textit{question}," which is, as we naturally expect, simply a state which is or is not unknown by a human, (being the mental analogue of an Aethic attribute, or perhaps the analogue of the type theory class). Then, there is an ``\textit{answer}," which is the state that said human assigns to the question, (being the mental analogue of an Aethic attribute's value, or perhaps the analogue of the type theory object). The key, here, is that we can suppose that a mind, when choosing an answer to a question, makes their choice based on the answers to other questions. For instance, ``Where are my car keys?" might borrow from the answer to the question ``Where did I last have them?" which itself borrows from ``Where did I have them last night?" and so on, until a satisfactory answer is found through this recursive process. This now being said, suppose we play this game of recursion, and wind up reaching a question for which the answer is firmly believed to be as such, but with no other question from which it borrows. Such a question is called a ``\textit{frozen question}," and its answer is then a ``\textit{frozen answer}." The classic example of its occurrence is during the ``why game" that a child might play. They progressively ask ``why" to every question's answer, until eventually you have no option but to reply ``I don't know." The question directly preceding this in the sequence, then, would be the frozen question, being the final question for which you have an answer, all while not itself holding a backup to its own answer beyond something like ``it just is." It is, in a sense, the foundation of all other questions which traced to it, but in the process borrows from nothing else -- hence triggering ``I don't know" when this is pressed by the questioner. 

More rigorously, a frozen question is intuitively known by a human to be true or in place, but by definition cannot have its answer yielded by any other question, answer, or piece of information that the human can readily access. As such, it has a connotation of being somewhat acceptingly unknown in derivation, metaphysically if not consciously -- but all the while firmly believed to hold some particular prescribed state. (Also note that to be a valid frozen question, it cannot be paradoxical or otherwise inconsequential to external logic, because such a thing would force it to be irrelevant to said human anyway). This being said, all we need to note at present is that, intuitively, a human mind holds at least one frozen question at any given moment -- or rather perhaps holds many of them. The second condition itself, then, can be stated in the following form: ``an agent may never hold less than one frozen question at a time." With the thought experiment, then, we can prove this logically. 

Let us imagine the question of whether or not there is a true objective morality to the universe. Either there is no such thing as moral objectivism, for which the question, ``what is the reason behind moral decision \textit{x}" immediately becomes a frozen question. Let us now show, then, why the other case, in which there is an objective morality to the universe, is also to imply at least one frozen question. In doing so, we will have that both possible cases result in at least one frozen question, hence proving our statement regarding such a thing. 

Suppose there is some objective morality to universe, as if it were a fundamental statement that was incorporated into the very foundation of reality. But let us ponder a simple query of such an objective moral truth: \textit{given that it is to exist, how might it be optimally communicated between agents?} If we assume that even a single agent, somehow, is in possession of this objective moral truth, (in a sense making them a sort of morally omniscient god), then is it possible in theory for such an agent to even communicate this truth to you or I? The answer, technically, is no. Even if they are in possession of all possible means of expressing the truth, we can never get over our own doubt that they might be deceiving us, because the only way to get around such a doubt would be to request yet another objective moral truth, which only they can provide us. We might imagine every possible means through which this moral deity expresses the objective moral truth to us -- perhaps they paint it across the sky, or perhaps they plant it into our dreams or even our subconscious minds. In all of these scenarios, however, we are fully unable to prove away any doubt of the statement being a deception. The furthest we can possible narrow it down to, is that it might be true, it might be deceptive, or it might be entirely irrelevant, (i.e. logically applying in some case but not others). The question of which of the three the stated objective moral truth is, then, is itself a frozen question, and to it we may only ever apply a frozen answer.

In effect, then, it is the combination of our being forced to make decisions of consequence together with our complete inability to get rid of our frozen questions that is the statement of agency's curse. 

Such things are a major problem I see in every philosophical argument I come across in person. Even the best philosophers one sees simply situate their frozen questions in results that they ``feel" are consistent, without ever rooting them to deeper ontology. As such, agency's curse looms over every discipline to which this type of procedure applies. No matter how much one believes their feeling is the right thing, so long as they do not root it to objective ontology, and their argument is negated by agency's curse. In effect, you get a room full of philosophers disagreeing about the specifics of every angle, which is definitely helpful for planning and sorting procedures, but is futile in the ontological lens due to agency's curse. And, it goes without saying that the breakage of the ontological lens is not a good thing. Applied over the scale of an entire civilization, such a thing can be catastrophic\footnote{And we have indeed seen it be so over the course of history, with all the wars and atrocities committed by humans. Notice how, every time, such things are a direct consequence of a group of people collectively agreeing on morally horrific ideals. However, even when we claim such a thing is morally horrific, we do so by referring to our own frozen questions, so in effect you get an infinitely spiralling system of nobody knowing if they are right, but all the while being sure, (or perhaps sure of being unsure in the case of absurdists), that they are. This will never move ground so long as agency's curse stands, because as we see in the Aethic paper \cite{Aethus}, progress is only as good as one's quality ontology. In effect we have no ammo to spin the debate in favor of one system of morality or meaning, and in the lens of objectivity we are no better or worse than the opposition. Like them, we have no access to objectivity as well. So who is worse, truly? (I mean, obviously the people who commit atrocities are worse, but the idea is that we can never back that up logically, and that's a real problem).} -- or not, because agency's curse does not endow you even with the knowledge of that. Specifically, this procedure becomes futile by making any perception of ontology no less arbitrary than random guessing. If you want to build a sustainable world in the art of ontology, it simply is not sufficient to have this kind of model. Hence the word ``curse." It is like driving utterly blind in our most intimate outlooks. I would argue that the presence of such an issue holds a significant piece of the existential burden outlined by the Postmodernist and absurdist movements. 

I would argue then, that agency's curse shares similarities with something like the Matrix, while in key ways being even worse. In a simulation scenario, with Aethic reasoning \cite{Aethus} we see that our Aethae are actually not demoted in their capacity of being ``real," and through this we would not suffer as much of a metaphysical blow as we like to assume -- with such a blow being both arbitrary and a construct. A more metaphysically concerning issue than this might be the problem of induction, which states that that which sits behind the entire array of induction done by human civilization may or may not operate as we can extrapolate from within. Still worse than this, perhaps, is the problem highlighted by epistemological skepticism, or its Aethic counterpart, in which the information held by the Aethus cannot ever truly be attained by a human with any sense of surety. In addition to this, we have the Aethic principle of accidentalism, through which metaphysics itself takes on the burden of these uncertainties in epistemology. But all of these issues, as unsettling as they may be, are not quite as looming over human existence as agency's curse, at least as is the perspective of this essay. Unlike these other problems, which at their worst can still be painlessly avoided or dealt with -- for example through us simply shrugging and reinterpreting our choices and or place in reality -- agency's curse is rather more Lovecraftian for its being like a parasite attached to any thinking being. This is why, at least in the epistemological sense, it can be viewed as the foremost philosophical concern of existing in the first place\footnote{To highlight this, we can even consider going through the generating thought experiment for agency's curse in just about any human discipline. We can create one for ethics, law, meaning, science, and so on, and while they are all variants of the larger curse, it does follow all the while that any single place where a conscious entity sticks their mind becomes corrupted by agency's curse when they do so.} -- because it lies beyond our interpretation of the outside world, and for that reason is more intimate with and deeper into the core of our respective consciousnesses. That is, it is about the metaphysical \textit{us}, not about the metaphysical \textit{outside} to us. 

It is quite acceptable at this point for a conscious being who has just philosophically engaged with agency's curse, (as we have just done), to spiral into a nihilistic dread where they state that life and death are collectively evil, however I would please ask that you refrain from doing so due to the next point -- being our saving grace. Yes! I am happy to inform that there is in fact a saving grace to agency's curse, if we so choose to use it. Of course, it is not an antidote to the curse, but it does serve as a perhaps miraculously sufficient countermeasure. If anything, it might have fallen under the radar if not for the curse itself, and as such is especially highlighted by it. The idea, then, is this -- given that we must select at least one frozen question, why not choose just such a question that our existence itself demands we give it a certain answer. How well can we do this? $90\%$ accuracy? $99\%$ accuracy? Can we form one single question which we are compelled to answer a certain way by a correlation through our existence, just as agency's curse itself is correlated through our existence? We have to think bigger than just opinions and core beliefs, because even those are fragile and separable from our content as entities. What is something that I physically cannot pull away from a conscious being? The purpose of the next section is to display just such a question -- \textit{The Golden Question}.

\subsection{The Golden Question}

What might the Golden Question be? And before you say that Decartes has already been over this when he said ``I think therefore I am," to that I say not by a long shot. First of all, Descartes was addressing whether or not he could determine if he exists. It's oddly specific and, Aethically, somewhat vague, biased, and unimportant. The Golden Question, then, is much broader to the workings of everything, and seeks to situate itself as one single question, technically being a frozen question, which we can root our entire system of morality and purpose around. Sure, there are some parallels to Descartes's procedure in gaining his proof of existence by answering all with one notion, but here, if anything, the stakes are far higher. 

This prelude being said, here is the statement of the Golden Question.
\begin{frozen}[The Golden Question]
Should our biosphere live forever?
\end{frozen}
It is that simple, but it can be seen as the sole point of transcendence beyond agency's curse. The reason for this is, unsurprisingly, mathematical. 

\subsubsection{Demonstration of the Relevance of the Golden Question}

The Golden Question is arguably the sole holder of what is needed to counteract agency's curse for one primary reason: no matter how we choose to answer it, our existence will, with time, tend toward us answering it in one definitive way, being the Golden Answer: simply \textit{yes}. 

To see the reason for this, we have to perform a mathematical derivation, which will itself yield insight into the mindset of our distant descendants regarding the way of things. The conclusion of this derivation can be stated in the form of what we will refer to as \textit{the Golden Tendency}.
\begin{prince}[The Golden Tendency]
Our descendants, should they exist, will tend with time toward existing in a society completely absent of the danger of their extinction, along with whatever consequences pair with such a thing.
\end{prince}
This can be considered to be a special case of the Lindy effect \cite{Taleb_Nassim_Nicholas2012-11-27}, but with the addition of further implied properties. For one, even if the hazard rates associated with these things are not precisely constant, (as they would be in the Lindy effect), and still this property will generally tend to apply, (although, by the Poisson limit theorem, hazard rates tend to be constant anyway, at least in any Aethus where such a thing is not implied to operate one way or another -- in essence making constant hazard rates the natural consequence of combining many random hazard rates together like white noise, as we see things do in reality due to complexity). Therefore, due to the intimate connection between this, and the overarching procedure of the Lindy effect, we see that the general Golden Tendency is applicable to any living system which exists over extremely long periods of time. For example, Tardigrades. The special form of the Golden Tendency, however, ought to apply to an entire biosphere, which itself operates over cosmic timescales rather than simply human timescales. 

In a sense, then, the mindset should be to draw a connection between it and our own human agency. We may reason that given this fundamental tendency, it follows that those biospheres which do tend toward having zero danger of extinction ought to be countering that with an organized conscious effort -- which itself is the direct ontological explanation for ``why" the human brain exists in the meaningful sense, at least according to Golden reasoning. This is a crucially important insight, because we can begin to use it as the basic element with which to construct a full system of meaning for our society. 

There is an insight of utmost importance that deduces the above notion for ``why" the human brain exists. To see this, consider any single stage of Earth's natural history, and then ask oneself the question of how such a stage ought to be transformed in the direction of compatibility with the Golden Tendency. For example, we might look back at the beginning of the Cenozoic Era, and ask where our biosphere at that time ought to tend over the next hundred million years so as to progressively minimize its own danger of extinction. Naturally, we might suppose that this biosphere ought to somehow generate a sub-component of itself which enables it to target and fight back against its own extinction pressures, such as the dangers present in interstellar space. Then, when we really look at that, we humans hold this precise property, given our agency and our brains. As such, the assertion of Golden reasoning is to, fundamentally, picture our biosphere not as some timeline from the distant past to the present, but rather as the natural geodesic of sorts between any given past iteration of our biosphere, and the optimal biosphere of the far future which abides by the Golden tendency. We will find, astonishingly, that our own biosphere will fall somewhere along this geodesic -- or at least that is a primary postulate in Golden reasoning. The idea, then, is to be taken in awe of the coincidence between the conditions met by the human brain and the conditions needed to get our biosphere to its far future goal. The fact that we coincidentally hold these conditions would be absurdly unlikely if the Cosmic Nexus were in control. As such, we are forced, by the mediocrity principle, to record a massive correlation between the Golden Tendency and the abilities of humankind itself. It is as if the state of Aetherians and the Golden Tendency are deeply connected below the surface. 

\begin{concept}[Golden Extrapolation]
    Ontologically speaking, we ought to picture ourselves as a node directly along the gradient from the Dimensional Null Aethus to the end-state of the Golden Tendency, rather than an isolated, culminating, or original group of entities. 
\end{concept}

With this statement of the importance of the Golden Tendency now being given, let us derive it, and then show how it pertains to the Golden Question itself.

\subsubsection{Derivation of the Golden Tendency}

Let us begin by noting a very intriguing mathematical concept, which will be generally referred to as \textit{the Golden Filter}. 
\begin{prince}[The Golden Filter]
    Any subset of the spacetime-Aethic combination that pertains to some child Aethus of an Aetherian will tend to, with time, be populated primarily by biospheres with an extremely high survival efficiency as opposed to not. 
\end{prince}
This concept has many sub-concepts, such as a version which applies to the future state of humanity, or another version which applies to the broad distribution of life in the universe. Generally, though, the procedure is always the same -- given some spacetime continuum, which is to be linked to an Aethus, it will follow that the number of biospheres in this spacetime which are efficient at self preservation will progressively tend with time to massively outnumber the rest of the biospheres. 

To intuitively see why this is the case, consider a fictional galaxy, for example, and for the purposes of the demonstration, we will tone down the Rare-Earth hypothesis to the scale of a galaxy. Suppose, then, that many different kinds of biospheres are to independently undergo abiogenesis far apart from one another, at different random spots in this galaxy. Over time, we will see that some of these biospheres tend toward being very unstable and having a low probability of continued survival. It follows, then, that these kinds of biospheres tend to, in all likelihood, fizzle out with time. More durable biospheres, however, will last longer, until they themselves fizzle out. For the sake of simplicity, if we imagine assigning a totally random hazard rate of death to every new biosphere, then it follows that the distribution with which these hazard rates is generated ought \textit{not} to match the distribution of abundance of these biospheres in the galaxy -- which itself will skew in favor of the more durable biospheres, which stick around for longer. We will notice more and more as time passes that the most durable of the biospheres tend to accumulate and stay, in a sense growing linearly in number, whereas the less durable ones will fizzle away in waves, essentially staying comparatively constant in number. And the same is true of the most durable even among the most durable, who themselves will display a similar pattern compared to the less durable portion of the highly durable. In general, we see a vastly propagating skew of the most durable biospheres becoming the most abundant in the universe, until, eventually, after enough time they essentially occupy the vast majority of life in the universe. This, then -- is the Golden Filter -- being the filter between the Nexus/distribution of generation for biospheres, (being the Cosmic Nexus), and the Nexus/distribution of abundance of biospheres, (being, at least in tendency, \textit{the Golden Nexus}, as we will call it). 

In the previous example, we have just applied the Golden Filter to the state of life in the entire cosmos, but, Aethically, we can do just the same for the Aethus of descendent biospheres to Earth's own biosphere. Rather than separating them through space in a galaxy, now we can separate them using a disagreeing superposition over their respective existences. Using this, let us now begin the official derivation of the Golden Tendency, and through that the Golden Question. 

Consider some Aethus which exists within and describes a version of our biosphere in the future, (so, technically speaking, consider that it is a child Aethus of one's own Aethus). Now, given this, we can seek to retrieve only two major values from this Aethus, with the first being the amount of time in the future for which we want said Aethus to situate, (as Aethae are timeless, suppose that we are considering some unit time to correspond to the Aethus, just as any current year might correspond to oneself), and the second being the hazard rate at said time of said Aethus's biosphere going extinct. Let us label these two values as the variables $t$ and $\lambda$. Then, given this, we might suppose, so as to conventionally mirror general centric unfolding, that any future Aethus is to be a child Aethus of some linear sequence of parent Aethae that correspond to previous times for $t$, and that $\lambda$ ought to be continuous over this sequence. Generally, then, note how no math has been created, and no stances taken in particular -- all we have done here is lay down a framework for rigorously describing these two particular attributes, which an Aethus must yield anyway through existing. And, we have taken advantage of the third postulate of the Aethus so as to suppose that said attributes operate in a disagreeing superposition with respect to their alternate states, that way the ontology of reality can be accurately mirrored. This now being said, let us suppose that we tentatively refer to the sequence of Aethae in time set up by this model as simply ``Aethic Lines." That is, every Aethus in this line will, by virtue of the child operation, imply the states of all Aethae previously based in the Aethic Line, but will hold all future possible Aethae in superposition. And, it ought to generate our $t$ and $\lambda$ values, which work to a specific default time over which that Aethus operates. (Regarding the Nexus which we use to operate these Aethic Lines, interestingly, it does not make a difference to upcoming final result whether we pick the Aetherian Nexus, the Cosmic Nexus, or whatever else within reason -- so that is up to the interpreter, but we will pick Dimensional for its ease of simplicity in the current derivation).

This now being said, let us set up a very specific probability distribution over the Aethus of all these Aethic Lines, (in effect being one's own Aethus that we started with, as it was defined that the Aethic Lines are themselves all possible child Aethae to said Aethus). First of all, let us consider some additional value to each one of these Aethae in each one of the lines, being some value of $s$. Suppose that we define $s$ as the ``survival probability" of the biosphere at said Aethus through its respective Aethic Line so far. As we know from survival analysis, $s$ can be written in terms of $t$ and $\lammbda$ as follows.
\begin{equation}
    s = e^{-\int_0^{t} \lambda \left( \tau \right) \, d\tau}
\end{equation}
Or, in terms of the Aethae in the line, we might write it as follows, in terms of the selected Aethus, being $A$.
\begin{equation}
    s = e^{-\int_0^{A_t} A_\lambda \left( \tau \right) \, d\tau}
\end{equation}
The distribution, then, is to, across all Aethae across all Aethic Lines which hold some time value of $t$, we might create the distribution of what the corresponding values of $s$ might be. That is, given $t$, we want to find the distribution of $s$ itself across all Aethae which are compatible with that $t$. Let us refer to a quantile within this distribution as some $x$. Using this, now, we can substitute $s$ itself with the notion of $x$, as after all through the definition of quantile, $x$ is a bijection of $s$ for some fixed $t$. As such, we can now write $s$ in terms of $t$ and $x$ rather than writing $x$ in terms of $t$ and $s$. Specifically, as a function of two variables, we get the following form.
\begin{equation}
    s = s \! \left( t, x \right)
\end{equation}

To intuitively understand $x$ at this point, consider it to be the ``efficiency" of survival of sorts over the path through scenarios in which a biosphere has taken. For some examples, if the biosphere were very unstable and constantly on the brink of nuclear disaster, then $x$ would be quite low. However, if it had organized a vast network of efficient self preservation instead, then $x$ would be respectively higher. The idea, then, is that in the ``set of all scenarios" for the future of our biosphere, $x$ tells where such a scenario's quantile is in the probability of the biosphere surviving. However, since such a sentiment is not entirely Aethically rigorous, the preceding definitions allow the classical intuition of such a phenomenon to be translated to Aethic terms. 

This all being said, now we have written $s$ in terms of the time in the future, $t$, and the efficiency quantile, $x$. Given such a trio of values, we might now also seek to find what the hazard rate is for such an Aethus, being $\lambda$. Note that there is no single value which it can take, because, plainly, many different sequences of $\lambda$ are possible to yield every one value of $s$. Nonetheless, we can find the distribution for $\lambda$, and through this we can find its expected value in terms of the $x$ and $t$ through which we calculate an $s$. Note that in the face of many possibilities over which $\lambda$ can function throughout time, itself representing an agreeing superposition\footnote{Agreeing and not disagreeing, because while the states they might describe are disagreeing typically, the $\lambda$ values are, through being probabilistic, therefore agreeing due to the third postulate of the Aethus} over them, we may approximate the presentation of such an agreeing superposition through the Poisson limit theorem itself -- or in other words, by taking advantage of the canceling of all that excess statistical noise of sorts. Through this, we tend to, as most things do in our universe, yield a presentation of that agreeing superposition which is roughly constant, thus yielding a roughly constant $\lambda$ to every Aethus in an Aethic Line, at least as far as the Poisson Limit theorem is concerned. %BACK UP THIS WITH MATH

This now being given, this approximation yields that a given $s$ value should imply a corresponding $\lambda$ value, specifically through the exponential distribution itself, with $\lambda = -\frac{\ln \left( s \right)}{t}$. As such, we can yield our new $\lambda$ in terms of the given function for $s$ in terms of $t$ and $x$.
\begin{equation}
    \lambda \! \left( t, x \right) = -\frac{\ln \! \left( s \! \left( t, x \right) \right)}{t}
\end{equation}
So, of course, for ease of notation, there is not much use in writing $s \! \left( t, x \right)$ now in the first place, as instead we can refer to $\lambda$ directly in terms of $\lambda \! \left( t, x \right)$. 

The question, now, is to see what we can decipher about $\lambda$ itself by virtue of any random processes that might ensue in humanity's future, themselves affecting $t$ and $x$. To begin to visualize such effects, let us attempt to consolidate the relation between $\lambda$ and $x$ by considering a Bayesian given: \textit{what if our biosphere is to definitely survive through time $t$?} With this in mind, we can consider that rather than $x$ itself uniformly random when describing a future by Aethic weight, instead we have added a new Nexus to the mix as well, being some ``Survival Nexus." And, intuitively enough, we might suppose that such a Nexus now rates the value $s \! \left( t, x \right)$ itself as the additionally valued weight of which futures now hold the most Aethae through this Nexus. As an example of such a Nexus in action, suppose we take the expected value of $X_t$ relative to this Nexus, where $X_t$ is a random variable representing the weighted QNumber of possibilities for $x$'s outcome. 
\begin{equation}
    \EX \left[ X_t \right] = \frac{\int_0^1 s \! \left( t, x \right) \, x \, dx}{\int_0^1 s \! \left( t, x \right) \, dx}
\end{equation}
\begin{equation}
    s \! \left( t, x \right) = e^{-\int_0^{t} \lambda \left( \tau, x \right) \, d\tau}
\end{equation}
\begin{equation}
    \EX \left[ X_t \right] = \frac{\int_0^1 e^{-\int_0^{t} \lambda \left( \tau, x \right) \, d\tau} \, x \, dx}{\int_0^1 e^{-\int_0^{t} \lambda \left( \tau, x \right) \, d\tau} \, dx}
\end{equation}
As an example of use for this function, consider an example possibility in which $\lambda \! \left( t, x \right) = {\left( 1 - x \right)}^{t}$, itself corresponding to an expected value of $x$ in the infinite future following $\EX \left[ X_\infty \right] \approx \frac{3}{4}$. Note how this is different from $\EX \left[ X_0 \right]$, which always equals $\frac{1}{2}$ on account of not having sufficient time to apply the Survival Nexus with any difference to whatever Nexus was set in place before, itself then setting a uniform random $x$ as per its quantile. Also note that upon creating a model in which the futures are collectively more dangerous, such as with $\lambda \! \left( t, x \right) = {\left( 1 - x^2 \right)}^{t}$, it then follows that $\EX \left[ X_\infty \right]$ takes the form of $\EX \left[ X_\infty \right] \approx \frac{5}{6}$ instead, hence arguing for only an increase in respective efficiency for those futures which survive. Note, however, that even this is something of an optimistic perspective for future survivability, because it assumes that all hazard rates across all $x$ values tend to zero with time, which is quite implausible in any semblance of realism.

Because of this, let us analyze one specific notion of such things. Consider some threshold of efficiency, $x_0$, which splits the Aethus of outcomes into two weighted categories: as either having less efficiency than it, or more. This may be represented by the following two integrals in the Survival Nexus, being $N_S$.
\begin{equation}
     P \left( X_t < x_0 \mid N_S \right) = \frac{\int_0^{x_0} e^{-\int_0^{t} \lambda \left( \tau, x \right) \, d\tau} \, dx}{\int_0^1 e^{-\int_0^{t} \lambda \left( \tau, x \right) \, d\tau} \, dx}
\end{equation}
\begin{equation}
     P \left( X_t > x_0 \mid N_S \right) = \frac{\int_{x_0}^1 e^{-\int_0^{t} \lambda \left( \tau, x \right) \, d\tau} \, dx}{\int_0^1 e^{-\int_0^{t} \lambda \left( \tau, x \right) \, d\tau} \, dx}
\end{equation}
Note that we can also write the Survival Nexus itself in terms of the background Nexus from which the Aethic Lines originally stemmed, that is given the specific condition that $S > t$, where $S$ represents a random variable of the time of death of a biosphere. That is, we may write in probabilistic notation that $N_S \equiv S > t$. Therefore, now consider the following set of statements.
\begin{equation}
    \frac{P \! \left( X_t < x_0 \mid N_S \right)}{P \! \left( X_t > x_0 \mid N_S \right)} = \frac{\int_0^{x_0} e^{-\int_0^{t} \lambda \left( \tau, x \right) \, d\tau} \, dx}{\int_{x_0}^1 e^{-\int_0^{t} \lambda \left( \tau, x \right) \, d\tau} \, dx}
\end{equation}
\begin{equation}
    \frac{P \! \left( N_S \right)}{P \! \left( X_t > x_0 \mid N_S \right)} = \frac{{\left(\int_0^1 e^{-\int_0^{t} \lambda \left( \tau, x \right) \, d\tau} \, dx \right)}^2}{\int_{x_0}^1 e^{-\int_0^{t} \lambda \left( \tau, x \right) \, d\tau} \, dx}
\end{equation}
\begin{equation}
    \frac{P \! \left( X_t < x_0 \mid N_S \right) \, P \! \left( N_S \right)}{{ P \left( X_t > x_0 \mid N_S \right)}^2} = \frac{\left( \int_0^{x_0} e^{-\int_0^{t} \lambda \left( \tau, x \right) \, d\tau} \, dx\right) {\left(\int_0^1 e^{-\int_0^{t} \lambda \left( \tau, x \right) \, d\tau} \, dx \right)}^2 }{{\left( \int_{x_0}^1 e^{-\int_0^{t} \lambda \left( \tau, x \right) \, d\tau} \, dx \right)}^2}
\end{equation}
\begin{equation}
    \frac{{P \! \left( N_S \right)} ^ 3}{{ P \! \left( X_t > x_0 \right)}^2} - \frac{{P \! \left( N_S \right)} ^ 2}{ P \! \left( X_t > x_0 \right)} = \frac{{\left(\int_0^1 e^{-\int_0^{t} \lambda \left( \tau, x \right) \, d\tau} \, dx \right)}^2}{{\left( \int_{x_0}^1 e^{-\int_0^{t} \lambda \left( \tau, x \right) \, d\tau} \, dx \right)}} - \frac{{\left(\int_0^1 e^{-\int_0^{t} \lambda \left( \tau, x \right) \, d\tau} \, dx \right)}^3}{{\left( \int_{x_0}^1 e^{-\int_0^{t} \lambda \left( \tau, x \right) \, d\tau} \, dx \right)}^2}
\end{equation}
\begin{equation}
    \frac{P \! \left( X_t < x_0 \mid N_S \right) \, P \! \left( N_S \right)}{{ P \left( X_t > x_0 \mid N_S \right)}^2} = \frac{ \int_0^1 \int_0^1 \int_0^{x_0} e^{-\int_0^{t} \left( \lambda \left( \tau, a \right) + \lambda \left( \tau, b \right) + \lambda \left( \tau, c \right) \right) \, d\tau} \, dc \, db \, da }
\end{equation}


\begin{equation}
    \exists k \! \in \R, \; 0 \leq \limsup_{t \to \infty} \frac{P \! \left( t > S \right) \, P \! \left( S > t > E^{-1} \right)}{{P \! \left( S > t \cap E^{-1} > t \right)}^2} < k
\end{equation}


\subsubsection{The Golden Question and the Meaning of Life}

The second major premise of Golden reasoning, now, is that the circumstance of the Golden tendency is not only part of an epistemic quandary, but also has a presence in a far greater struggle which has been unfolding since the early days of the universe. This is the second leg of Nietzsche’s serpent, which we will refer to as the \textit{Departric tendency}. 
A fundamental premise of Golden reasoning is that biospheres are not like humans, who always know they will die, and seek to live life to the fullest in the meantime. In Golden reasoning, we characterize a biosphere which will certainly inevitably die out as no less dead than a biosphere which has physically died already. Such an assessment simply follows from the indiscriminate time principle of Aethic reasoning, where the knowledge of a biosphere's death is equally valid no matter whether it comes before or after the physical event itself. 

% Fit in with this other content...

Now the Golden Tendency can be logically displayed -- and through this we see that our distant descendants will, as a consequence of their existence, tend strictly toward lowering their own danger of extinction to zero through their own conscious efforts. This notion yields a myriad of questions, many among them being philosophical. For instance, we might ask if these descendants have free will at all, if the nature of their beliefs can be so deterministically summarized by myself in the 21\textsuperscript{st} century at a laptop. This is a very interesting question, and it seems that it will yield only clearer answers with time, but at the moment my only instinct is the suppose that perhaps our modern notion of free will itself is archaic compared to some of the philosophy that our distant descendants in question will be grappling with. So, if there were same way to tailor it to their system, we would be wiser to assume such a way exists and is not known to us, than does not exist due to not being known to us. 

Another major point of contention that I see with the Golden Tendency, is that response which almost everyone I have spoken to about it reacts with -- \textit{humans are naturally chaotic, and many possibilities are open, so something this deterministic should not be known here and now}. To respond to this notion, yes, I understand that view, but it is actually part of the core ontological intrigue of the Golden Tendency that it can bypass such an intuition. See, the Golden Tendency is fundamentally a mathematical statement, concerned with order, but that doesn't mean it cannot be brought about by chaos, being the sum total of all human agency. History has told us that such a thing is abundantly possible elsewhere, for example the ideal gas law -- something entirely brought about by the random motions of molecules -- and yet something which, when all gathered together, yields an intriguing order to it. The precise same kind of thing, then, is what is true of our entire biosphere as regards the Golden Tendency. And, by extension of this, I would argue that a major shortcoming of human philosophy at this point in history is to focus so much on the individual molecules, that we lose the nature of the ideal gas law happening all along. So humans may very well be chaotic, and we may very well possess an abundance of free will -- but it stays the same that, mathematically, as sure as the sun is to rise in the East -- the Golden Tendency is to shepherd the future of our biosphere. 

And this brings us, of course, to the Golden Question itself. Consider the following dichotomy -- either our biosphere is to survive, or it is to go extinct, all by some arbitrary point of checking in the future. If our biosphere is to go extinct, then the meaning to its further existence is as clear as that -- none at all. So that makes one half of the dichotomy explained. And we have now also shown the nature of the other half of the dichotomy -- being survival -- through the use of the Golden Tendency itself. There is now something of a key philosophical insight to undergo at this stage. Take the Golden Question itself: ``\textit{Should our biosphere live forever?}" Note, intriguingly, how every single state of our biosphere which even has the capacity to answer such a question, in possessing such, rests on the ``survival" half of the dichotomy. Nexusiveally, one will note that every single potential future in which we exist, hence deciding whether to comply with the Golden Answer, will follow the Survival Nexus rather than any kind of broad Nexus over the entire dichotomy. Or, more simply, into every single system in which we place the asking of this question, we must place ourselves into such a system as well, to perform the actual doing of the ask. We are, then, fundamentally inseparable from the survival side of the dichotomy, because the state of us being relevant in the first place implies only the relevance of the survival half of the dichotomy. Through this, then -- we see that the Golden Tendency is perpetually relevant to us, because any us where there is, is a place where the Golden Tendency operates. Through this line of reasoning, then, we can combine the clear and unshakable tautology of the Golden Tendency with the clear and unshakable relevance of ourselves with it, and in the process, imply a definitive asserting of \textit{yes} to the Golden Question. In essence, the very state of life in the universe must, without exception, tend toward answering it this way, and so it is the single question in the entire universe which is utterly deterministically static in life. It is, in effect, the trajectory, and as such the meaning\footnote{That is, at least according to Golden reasoning!} of life to reach a state where the Golden Question can not only be answered with \textit{yes}, but be acted upon as such. 

To be abundantly clear, this is presenting itself not as some substitute or stand-in of sorts to life's meaning, but as the answer we have been searching for all along. That is, our metaphysical presence here, on Earth, is regarded as being directly associated with a long ray of life's endurance into the far future, rather than being some meaningless or trivial phenomenon. That is, instead of keeping the epistemological slot labeled ``meaning" as blank or unimportant, we can directly place into it this notion of trajectory towards a state in which life has essentially stagnated its own entropic dissolution, which we all contribute to just by existing. The vast majority of attempts\footnote{Meaning biospheres.} to reach this milestone end in failure and utter destruction through the Cosmic Nexus, however some few approach it with utmost closeness. In a sense, then, we, here and now, are that vitally important intermediate stage in this process, where life transitions to this higher ability. But it is truly a deterministic-leaning struggle, as old as time itself, which our existence within is the basis of our metaphysical content. 

Regarding agency's curse itself, we can view the Golden Question as the single static point which is both a frozen question and a question with a sure answer, at least as per our own existence. This, arguably, can be considered the optimal countermeasure against agency's curse, and so, somewhat fittingly, can be the avenue through which we rise above our current position in the broad spectrum of agency, with nonliving matter at the root, and omniscience itself at the peak. It is widely supposed by futurists that we are headed somewhat this way, and perhaps the transition from an agency's curse era of life to a Golden era of life can serve as the driving force behind just such a transition, both metaphysically and epistemologically. 

So, given this concept that our future, should it exist, ought to hold the Golden Answer in utmost prestige, we can use such a notion to gain insight into the lives of our distant descendants. 

\subsubsection{Alternative Phrasings of the Golden Answer}

\subsubsubsection{A Philosophical Scientific Method}

I was once asked a question of how we know the scientific method itself is not ascientific. That is, considering how that which is science is founded on its own set of principles, it follows that we might be able to create the same arguments against it that we may for ascientific concepts. Let me explain, however, what the scientific response would be to such a question, but put into philosophical terms: in claiming that we hold the inherent capacity as agents to possess and transfer information, we are implying that said information and its corresponding knowledge must be directly compatible with the scientific method. If we were to claim the alternative, then we would be creating an Aethic contradiction in our very supposition that we are allowed, as agents, to possess knowledge at all. Therefore, we must answer that science is not ascientific in the affirmative, because otherwise we are giving rise to a paradox. More specifically still, though, the reason that knowledge is dependent upon the scientific method is that knowledge is gathered from direct empirical observations, so for us to create a framework in which the knowledge and the empirical observations are distinct is for us to create a framework in which our knowledge it utterly uncorrelated to reality, and as such is nonexistent with respect to reality and through such things, ourselves. However, of course, even supposing that empiricism is empirically based on empiricism is implying that the scientific method aligns with empiricism, which itself relies on the conclusion as part of the derivation if we look at it through that lens. In effect, then, the framework we have created in supposing that science is not ascientific is circular logic which, by its own measure, is the least circular of all logic. Considering how no single other known argument holds such a property besides from the scientific method, (as under all alternative known arguments, all such arguments are evenly implicative and in agreement with themselves as in disagreement), we can therefore conclude that skepticism still reigns supreme, but the scientific method is just barely enough ammo to semi-combat it. And, from this line of logic, (subconsciously if not deliberately), we have founded the entirety of modern civilization. What I want to suppose, then, is that the corresponding revolution in the field of philosophy is actually yet to have been undergone, because where science created the scientific method to combat skepticism, philosophy never created a system to combat agency’s curse. That is, where agency’s curse is the philosophical analogue of skepticism, there is as of yet no philosophical analogue of the scientific method. But this is what makes Golden reasoning so exciting – because it has the sole potential to take on this role. 


\section{An Analysis of Future Tendencies}

\subsection{The Golden Initiative}

It is somewhat simple to see how in the far future, once our descendants have drawn their own hazard rates of extinction to near-zero, it follows that they will hold a strong initiative toward keeping the biosphere alive. Let this be stated. 
\begin{concept}[The Golden Initiative]
    The initiative of our biosphere in the future towards aligning practical outcomes with the abstract philosophy of the Golden Answer. 
\end{concept}
This now being said, as is implied by the low hazard rate, it follows that the Golden Initiative will tend to only ever grow stronger with time, until, by the far future, it is essentially the central component of our entire society, perhaps serving as the analogue to the modern ``self preservation" desire. 

It therefore follows that the Golden Initiative serves as the conscious-agency component to what the Golden Tendency describes mathematically.

\subsection{Golden Biospheres}

Another key term we can use in the description of the Golden Tendency, is what we will refer to as a ``Golden Biosphere," which is described as that biosphere to which the Golden Tendency tends a biosphere toward. Due to the nature of infinities being difficult to describe rigorously, it is most useful in the practical sense to consider biospheres with progressively lower hazard rates of death as, in a sense, ``approaching" the state of being a Golden Biosphere. That is, in this way we can use it as a mode of comparison, arguing that biosphere $A$ might be ``closer to" being a Golden Biosphere than biosphere $B$, as per some set of conditions, (primarily being the Aurutance of $A$ versus $B$, as will be described later on). 
\begin{concept}[Proper Golden Biosphere]
    The equivalence class of biospheres which operate at the fixed point at the end of the Golden Tendency, where the limit supremum of hazard rates of destruction are zero. If not a direct description of such biospheres, it is mostly that which other biospheres are described as ``tending toward" through the Golden Tendency. 
\end{concept}

This is considered to be the strict definition of a Golden Biosphere -- that biosphere which all extremely efficient biospheres tend toward and can be compared against. Regarding a looser definition, then, we might consider the \textit{Improper Golden Biosphere}, also sometimes just called a "Golden Biosphere" anyway, which is considered to be any biosphere which is of phenomenonally extreme levels of self-preservation, in a sense being within a very small margin of error away from a Proper Golden Biosphere. Of course, this is somewhat subjective due to the notion of the choice of a margin of error, however generally this term can be applied to biospheres which are conceptually indistinguishable from a Proper Golden Biosphere as per the descriptive system of choice. Unlike Proper Golden Biospheres, which represent a mathematical extreme, Improper Golden Biospheres are actually practically achievable, and as such represent the physical manifestation of the concept of a Golden Biosphere, (just as a planet might represent the physical manifestation of a sphere, even while not perfectly being a sphere). 

Note, then, the concept of a \textit{Golden Point}, which is mostly a descriptive term, which represents the point in Aethic time at which a biosphere passes the threshold required to be an Improper Golden Biosphere. Through this, if we assume that an Improper Golden Biosphere's approximation of a Proper Golden Biosphere is a tautology, (similarly to how we treated the mediocrity principle in Nexusive reasoning), then the following approximate definition for a Golden Point can be written. 
\begin{concept}[Golden Point]
    The point after which a biosphere is too capable at self-preservation to regress past that point.
\end{concept}
In essence then, we suppose that a fundamental criteria for a biosphere becoming a Golden Biosphere is that it ought not be able to regress past that point again. (In essence, then, a Golden Biosphere implies the state of being a Golden Biosphere indefinitely onward from there). For an Improper Golden Biosphere, we can take this to mean that said biosphere will never regress out of being an Improper Golden Biosphere for the foreseeable far future. 

\subsubsection{Comparison to the Technological Singularity}

I think that the Golden Point and the technological singularity \cite{} certainly have some major similarities, but there are also two primary differences in my view.
\begin{enumerate}
    \item The technological singularity is often framed in the context of the period between revolutions shrinking exponentially, until we hit a single asymptotic limiting point, (ie rather like the Feigenbaum point in the bifurcation diagram of the logistic map), and I do indeed think this is a shared and central characteristic of the Golden Point as well. However, my main issue with the technological singularity at this moment is that it stops just short of then taking the major contextual leap into framing this emergent point. That is, the standing interpretation is simply ``complexity", but this feels uncomfortably dry to me for what should really be the epitome of philosophical richness in practice. At the very least, I recommend asserting a semantic axiom for the ``ontological quality" of this event which captures it in a single abstract object, but better still is to then translate this axiom into an extension of some present-day context. 
    
    My proposal for just how to do this is the statement of the Golden initiative itself, as was inductively motivated by the Nexusive-Golden crossover: that is, that we are not at the beginning nor the end of the biological circumstance, but rather lie upon a predefined trajectory for which a Golden biosphere is optimally constructed by the universe. So, strictly speaking, the Golden Point is not a Feigenbaum point for biological complexity alone---rather it is the establishment of cosmic survivability, set in stone to a Nexusive degree of scope, which merely presents as a Feigenbaum point from the perspective of those ``within the Aethus". So, similarly to the fundamental theorem of empirical Aethae from Aethic reasoning \cite{Aethus}, at the highest ontological layer the Golden point is merely an optimization point in a vast ontological static Aethus of digits and intangible structures---but upon sparking ``experience" into this Aethic structure, one becomes an agent in the direct preceding events who views a Feigenbaum point on the horizon. So, if anything, the physical manifestation of the Golden Point is its peripheral appearance rather than its central Aethic structure. But the absolute key of understanding its richness lies in viewing it as a contextually immortal optimization point in a field of Aethic algebraic complexity, for which the Nexusive accordance principle draws us to its vicinity, and the Aethic translation to realization presents it to us as a Feigenbaum point---all while none of these circumstances even begin to capture its intrinsic depth of nature. In that way, it is something like attempting to describe all the ocean via an exposure to one sea shell: one is a subjective lens on a single product, and the other an abstract-ontic generative superpower. 

    We might even refer to this ontological priority of the Golden Point as something of a principle.
    \begin{prince}[The Golden Well Principle]
        Fundamentally speaking, the Golden Point is by far the most overwhelming and ontically primordial object in our empirical Aethic structure, rather analogously to how the Sun is our central gravitational beaco. In this sense, we might refer to it as our ``primary ontic well", Aethically speaking, in the sense that it is of a vastness and depth that we may infer as beyond anything else in our Aethic sphere of influence, (and even gives rise to that sphere of influence), but all the while is somewhat mysterious beyond this. At center, this feature of the Golden Point is localized directly to its state of fundamentally designating how our empirical static Aethus of the universe for some reason appears to be arranged around the Golden Point, such that during the static Aethus-to-conscious experience translation the Golden Point appears as an empirical Feigenbaum point upon existence. So while this Aethic object has a depth beyond what may be fathomed, the context around which it came into existence, or among what external world it exists, remains obscure\footnote{For now, at least.}. But if we were microbes drifting in a metaphorical dewdrop universe, the Golden Point---as a parent Aethus---is the very dewdrop. 

        Indeed, to directly play into the gravitational analogy, we have that trying to imagine the Solar System without the gravitational field of the Sun is rather analogous to attempting to imagine our Nexusive empirical context without the Golden Point. The major supposition is that the cohesion of things would crumble, and we would no longer have any Nexusive priority toward existing here at all\footnote{So altogether, if we imagine treater the interior structure of our own empirical Aethae like a ``field" of sorts, the Golden Point orients this field to the same degree that spontaneous symmetry breaking orients the empirical Higgs field: that is, by introducing this extra dynamical layer to our Aethus may we spark the cascade which leads to consciousness, but never without it. The Aethic postulates would still exist, but they would be forced to run on an unfeeling uniform soup rather than a structurally rich conscious landscape. But yes, altogether the Golden Point is the proposed root of this mechanism, however in particular it is contextually able to achieve this.}. Then, the Golden Well Principle itself is the further supposition that the Golden Point as an Aethus is almost like a mysterious and obscure artifact of unknown origin, (perhaps similarly to the monoliths in \textit{2001: A Space Odyssey} \cite{2001spaceodyssey}---with the major difference being that its causality points backwards from the Aethic future rather than forwards from the temporal past), in the sense that it is perhaps far and away the single most clearly ``primordial" or at least ``raw" metaphysical structure we are currently aware of. Like something made of a more ancestral substance than anything which has existed since, if you will. 
    \end{prince}

    \item Arising directly from this distinction with the technological singularity, then, is the second major difference of the Golden Point, \textit{which is that it now provides us with a sweeping framing of the present day, due in particular to its being ontologically deeper than anything else in our Aethus}. That is, rather than us being in an accidental materialist universe, with the technological singularity apparently happening to lie in our future, instead the contextual framing reverses: the reason we are Nexusiveally compelled to exist here in particular is \textit{because} the Golden Point lies in our immediate future, and so from that perspective it serves as the immediate node from which we are generatively reverse-engineered. But, even beyond the ontological sector, we also have the ideational sector as well, for which the Golden Point allows for us to frame existing circumstances in the context of how they retroactively interact with it. So, rather than the technological singularity being a passive peculiarity, the Golden Point is both fertile and the single most active player in the current circumstantial structure of reality, (as per its Aethic context). In this respect its ascension to relevance is absolute, even and especially now. 
\end{enumerate}

\subsection{The Departure Point}

With Golden Biospheres having been described, let us analyze the opposite of this effect -- being a biosphere which has already died out rather than having eliminated the chance of its own death. There are a few immediate concepts which come to mind to describe such a thing, such as an apocalypse, or a generally biosphere-wide extinction. However, note that these terms are indeed somewhat vague, especially with regards to their association with a Golden Biosphere, so perhaps a more general and explanatory term would be more preferable. Consider the concept of a \textit{Departure Point}, which, as the name demonstrates, constitutes the action of a biosphere departing from any potential to become a Golden Biosphere. A Departure Point, then, can be thought of as an event horizon which stands between a biosphere and the Golden Point. Let us describe this as such. 

\begin{concept}[The Departure Point]
    The point after which a biosphere is too incapable at self-preservation to progress past that point. 
\end{concept}

\begin{concept}[Departure Biosphere]
    A biosphere which has crossed the Departure Point.
\end{concept}

Regarding the nature of the word ``progress" in this context, it simply means the ability of a biosphere to increase its own Aurutance, (where Aurutance, as will be described in further detail later on, is a linear measure of how close an Aethus tends to implying a Golden Biosphere, with 0 meaning no effect, negative meaning a decrease in the capacity, and positive meaning an increase in the capacity). Like how the Golden Point, as applied to Improper Golden Biospheres, merely represents a very low margin of error from a Golden Biosphere, we can also suppose that the Departure Point is a fundamentally approximate object. Nonetheless, however, it is not difficult conceptually to identify or label a practical circumstance with a Departure Point. For instance, if we were to analyze the derivative of a biosphere's Aurutance with respect to Aethic time, we might suppose that a catastrophic drop-off in the Aurutance over an extremely small period of time serves as a Departure Point -- so perhaps we would conventionally situate the Departure Point itself at the inflection point of such an Aurutance function of Aethic time. Note, however, that the use of a Departure Point in a system is not to be taken lightly, as this can be seen to represent the highest form of disaster for a biosphere. As such, it should be reserved for those systems which are essentially irreversibly destroyed, putting any chances of recovery on the order of trillions upon trillions of times lower than standing optimal probabilities of biosphere immortality. Note that, by the analyses seen in the Nexusive reasoning paper \cite{Nexus}, such a thing can happen in a flash, with a droppage of indefinite survival probabilities by dozens of orders of magnitude being extremely possible under the Cosmic Nexus. We simply refer to such a thing as an \textit{Imperfection Magnitude Drop}. A Departure Point, then, can be considered to be the strongest tier of Imperfection Magnitude Drop, and therefore can be described as such. 

Now let us consider the philosophical context of these things.
The Departure Point, see, is not the destruction of life on Earth: \textit{it is the inevitability of the destruction of life on Earth}. At whichever point this inevitability first takes effect, then, is where it would lie — not with the destruction event itself. The Golden Point itself, then, can under this lens be visualized as the point after which a biosphere can never again backtrack in the direction of the Departure Point, in effect rendering not only the Departure Point, but the very possibility of approaching the Departure Point as an implausibility. Guaranteeing this to practical certainty is what marks the Golden Point. 

Next, please note the existence of the following theorem, with it merely being a consequence of the mean value theorem. 
\begin{theorem}[Departric Dichotomy Theorem]
    Any given biosphere which fails to ever reach its Golden Point must inevitably cross a Departure Point, (with the more trivial converse also being in effect). 
\end{theorem}
Regarding the specific proof of this, for a biosphere to never reach its Golden Point means that its probability of becoming a Golden Biosphere is bounded from above by some $k < 1$. 
\begin{equation}
    P \! \left( G \mid T = t \right) < k < 1
\end{equation}
This is so for some Golden Aethus of $G$, and an Aethic time of $T$. We further have the following property on account of the Golden Filter principle.
\begin{equation}
    \lim_{t \to \infty} P \! \left( G \mid G_T > t \right) = 0
\end{equation}
This is to say that if the Aethic time at which a biosphere reaches the Golden Point, if at all, is some $G_T$, then if this has not already happened by $t$, we have that the probability of it happening anytime after $t$ goes to zero as $t$ goes to infinity. We can show the validity of this as simply as noticing how its contrapositive is essentially equivalent to the statement of the Golden Filter itself, where the probability of infinitely surviving biospheres being Golden goes to one hundred percent. 

As such, we have that a biosphere's probability of becoming Golden with respect to some Aethic time $t$ is both bounded from above by $k < 1$ and destined to drop to zero if said biosphere never reaches a Golden Point. Hence, if we pick an arbitrary Aethic time $t$ to analyze, then at such a time the probability of the Golden Point might be some $P \! \left( G \mid T = t \right)$, which we know follows $0 < P \! \left( G \mid T = t \right) < k < 1$. At this point we have two options: either the probability of the Golden Point will never again rise above $P \! \left( G \mid T = t \right)$, which would make $t$ itself a Departure Point by definition, or instead the probability \textit{will indeed} rise above $P \! \left( G \mid T = t \right)$ at some point in the future. In this second case, we have by the intermediate value theorem that there exists a future Aethic time of $t_2$ for which the probability exactly matches what it was at $t$. This is simply because the probability would have risen after $t$ as per the second case in question, and then after that must have lowered so as to tend to zero. The intermediate value theorem itself then implies that the intermediate probability of $P \! \left( G \mid T = t \right)$ must lie between the two boundaries.
\begin{equation}
    P \! \left( G \mid T = t \right) = P \! \left( G \mid T = t_2 > t \right)
\end{equation}
Finally, we get from the mean value theorem that between $T=t$ and $T=t_2$, there must be some $t_1$ for which the derivative of the probability with respect to $T$ is zero. We also know that at least one such possible $t_1$ must have probability greater than that of $t$, because of how the derivative of the probability itself was positive at $t$, and then negative at at least one $t_3 = t \to \infty$. Hence, in this case two we have a definite $t_1$ for which the following holds. 
\begin{equation}
    P \! \left( G \mid T = t \right) < P \! \left( G \mid t_2 > T = t_1 > t \right) < k < 1
\end{equation}
As such, if we consider the set of all choices for $t_2$ which can be derived with the mean value theorem in the same manner, then the supremum of this set is itself guaranteed to be greater than or equal to all remaining relative maxima, so by the transitive property is implied to be the global maximum of the entire probability function of Aethic time. As such, for any $t_0 > t_1$, we have the following.
\begin{equation}
    t_0 > t_1 \Leftrightarrow P \! \left( G \mid T = t_0 \right) < P \! \left( G \mid T = t_1 \right) < k < 1
\end{equation}
As such, $t_1$ meets the definition of a Departure Point.

We see with the Departric dichotomy theorem that there is a logical dichotomy between a biosphere becoming a Golden Biosphere versus becoming a Departure Biosphere. That is, they are mutually exclusive and collectively exhaustive outcomes. This makes the specific relevance of the Departure Point all the more apparent, because it is quite literally -- to a mathematical degree of rigor -- the antithesis of the Golden Point. These are among the reasons why we argue that it would be beneficial to the human paradigm for our fundamental dichotomy of biosphere-existence to be \textit{Golden versus Departric} rather than \textit{extant versus extinct} or the like. It is this connotation of event horizons that yields the most ontological insight into the workings of such phenomena, and as such is likeliest to be adopted by our descendants under the Golden Tendency. 

\subsubsection{Statement of the Golden Imperative}

Given the relational definition between the Golden Point and the Departure Point, as well as the dichotomous stance of exactly one of them being given to happen in the future, we might consider the phrasing the importance of the Golden Point with respect to the capacity of the Departure Point to happen.
\begin{prince}[The Golden Imperative]
    The only way to prevent the Departure Point is to reach the Golden Point. As such, we might picture life itself, in cosmic terms, as representative of the immortal struggle between the Golden Point and the Departure Point. As we know from the Rare Earth hypothesis, the Departure Point wins almost every time, so the Golden Point itself can be viewed as the infinitely sought after immunity device against the Departure Point. 
\end{prince}

\subsubsection{Specific Interpretations of the Golden Point}

Using this philosophical definitions of the Golden Point and the Departure Point, there are a few possible mathematically well-defined candidates which amount to a definition of the Golden Point. Let us consider these.
\begin{itemize}
    \item Philosophical Golden Point
    \begin{itemize}
        \item \textit{The point after which a biosphere is too capable at self-preservation to regress past that point.}
        \item This is the base template of the Golden Point's definition. 
    \end{itemize}
    \item High Golden Point
    \begin{itemize}
        \item \textit{The first point in Aethic time at which a biosphere is a Proper Golden Biosphere.}
        \item This is somewhat impossible to physically happen, but it is nonetheless useful as a mathematical construct when used interchangeably with the state of a Proper Golden Biosphere existing. (For example, we might say that the probability of a Proper Golden Biosphere existing equals the probability of reaching the Heavy Golden Point).
    \end{itemize}
    \item Light Golden Point
    \begin{itemize}
        \item \textit{The first point in Aethic time at which a biosphere is an Improper Golden Biosphere}.
        \item Through being defined in terms of an Improper Golden Biosphere, the Light Golden Point inherits the subjectivity in its definition. Generally, this is meant for the allowance of a subjectively defined margin of error in the resemblance between one's own biosphere and a Proper Golden Biosphere. 
    \end{itemize}
    \item Deterministic Golden Point
    \begin{itemize}
        \item \textit{The first point in Aethic time for which all future Aethic times hold a higher probability of reaching the High Golden Point than it does.}
        \item The benefit of this approach is that it incorporates the event horizon-like connotation of a Golden Point, however it does not quite capture the agent-based capability associated with the Golden Point. That is, if we consider the two notions of self-preservation and progress versus regress in the philosophical definition, then to define regress as a biosphere decreasing its probability of the High Golden Point is a very different story than defining regress as a biosphere decreasing its own instantaneous ability to prevent a probability decrease. The latter seems more on par with the notion of self-preservation, being a measure of one's own prevention of threats more so than a measure of the mere frequency of realized threats. As such, the Deterministic Golden Point misses the mark on this on account of instead associating with the former. 
        \item It should further be noted that this approach is not compatible with Aethic reasoning in the first place. That is, it relies on a deterministic and objective nature of data which is not sufficiently Aethic. For such reasons a Deterministic Golden Point will only ever be approximate in scope. 
    \end{itemize}
    \item Proactive Golden Point
    \begin{itemize}
        \item \textit{The first point in Aethic time for which all future Aethic times hold a lower Departric hazard rate than it does.}
        \item Even though this is still a bit too deterministic for rigorous Aethic reasoning, it is a marked improvement over the standard Deterministic Golden Point definition for its closeness to the philosophical definition of the Golden Point. If we choose to measure the ability of a biosphere to engage in self-preservation as being inversely proportional to the hazard rate of the Departure Point occurring, then for a biosphere to regress in self-preservation would mean an increase in this hazard rate. As such, if we suppose that the hazard rate is to never again exceed the value that it holds at a particular Aethic time, then we do indeed have that that Aethic time represents a point at which regression in self-preservation is permanently prevented. 
    \end{itemize}
    \item Expected Golden Point
    \begin{itemize}
        \item \textit{The first point in Aethic time for which the probability of its Departric hazard rate ever being exceeded again is less than half.}
        \label{expect Gold}
        \item This definition serves to take the basic principle of the Proactive Golden Point, and then make it compatible with Aethic reasoning. Rather than deterministically asserting that the hazard rate function itself for a biosphere must hold a fixed value, we are instead treating such a hazard rate function as a random function. Further suppositions about such a stochastic process will be made during the \hyperref[stochastic analysis]{future tendency section}.
    \end{itemize}
\end{itemize}

\subsection{The Broad Landscape of Life in the Universe}

With the mathematical position of the Departure Point described, now we must analyze it through the tangible lens. A classical way to think of the Departure Point might be to consider it something recently relevant — as a kind of dark finish line we can either cross or avoid due to human ingenuity alone, (as if it were fully synonymous with a human-generated apocalypse). But, this characterization is mistaken. The truth is, we have been at war with the Departure Point for the past several billion years now. It is the oldest and most primordial enemy of life in the universe. 

By Golden Reasoning, the outlook of the cosmos is essentially this: life, from its earliest origins on a planet, undergoes a race to reach its Golden Point before the Departure Point claims it. This gives needed context to the actual ontology of what life \textit{is}. Life, through Golden Reasoning, is a system which is in disagreeing superposition between undergoing a Departure Point in the future, or a Golden Point instead. If it reaches a Departure Point first, then we declare it, ``extinct," whereas if it reaches a Golden Point, we now designate the title of ``Golden Life" to it, which can be considered the next stage, or perhaps end state of life. But so long as a life-form remains within this superposition, it is only in an intermediate stage between inorganic matter and Golden Life, much like how we humans might describe a fish as roughly an intermediate stage between a Prokaryote and a human being. Given that the vast majority of life in the future cosmos is to be Golden Life, it then follows that we ourselves are not actually ``complete" in any sense of the word, either in surety of remaining in existence, or perhaps even in the taxonomic sense, where we share many similarities both with inorganic matter and Golden Life, without falling in either category. But the essential notion to realize of the future universe is that the fundamental dichotomy of things is between abiotic material and Golden Life -- not between abiotic and biotic material, with such a thing being comparatively arbitrary because it is simply not an extremum by any measure\footnote{Or at least, not by any fundamental measure. Due to chaos, one could make an extremum out of anything by choosing the right function, but of course some functions are more arbitrary than others. Functions like abundance in the universe, hazard rates of death, or evolutionary progression are all examples of leaning toward fundamentality.}. So, compared with Golden Life, we ourselves at this moment are like the crude algae one might find in a cave underneath the surface of the Earth -- locked away in some quiet corner of the universe, and deeply dependent on our environments and the general luck of the draw. At the very least, our descendants might regard us like how we regard the hypothetical RNA world \cite{} -- a stepping stone, although deeply impactful to all which came after it. And for this reason it would not be quite accurate to call ourselves the most important generation in Earth's history, because in a way every generation and every achievement was equal in its importance to the Golden Point. However, we still may very well be the most important generation in \textit{human history}, because we are among the first that will mindfully contend with some of the stakes that will someday be signature of Golden Life. 

More on this point, let us briefly analyze what it is that we have to contend with in the first place.

\subsection{Internal and External Dangers}

As was briefly elaborated upon in the paper on Nexusive reasoning \cite{Nexus}, the Departure Point itself can be subdivided into two distinct categories, being the internal Departure Point and the external Departure Point. The differences between the two are just as the names describe -- for internal Departure Points, the source is within the biosphere, with some common examples being human induced apocalypses, or perhaps even general evolutionary stagnation, (with a connotation of some disaster occurring in the process of that). As for external Departure Points, they come from beyond the biosphere, and inflict death upon it from space. Examples of this could be asteroid impacts, supernovae, and so on -- but at the end of the day, the deadliest external departure points are the most mundane -- the orbital or habitability destabilization of the planet. There are a trillion possible ways things could go wrong at any given moment. Right now, for instance, the Solar System is moving through a hailstorm of other stars which are zipping around next to us. Over geologic timescales, it is something of a miracle that at least the orbit of Neptune has not been thrown off due to the gravity of passing stars, taking the whole rest of the Solar System with it. Generally, we see that Earth probably sits at an unstable equilibrium point rather than a stable one, (just due to the logic of what was outlined in the Nexusive reasoning paper), and we can argue that the external Departure danger is the very root of such an issue. 

These terms now being defined, let us demonstrate how they are relevant to the state of the Golden Initiative. To begin with, it follows as a lemma of the Golden Tendency that any Golden Biosphere ought to have figured out a way to fully nullify both, as if there was any remaining danger in either category, such a biosphere would not qualify as a Golden Biosphere in the first place. As such, one of the primary tools we can use to characterize a Golden Biosphere for our research as modern humans, it can be argued, is to ask what a biosphere looks like which has nullified both internal and external dangers of its own destruction. 

The first thing to note about this is that there is something of an odd discrepancy between what lowers internal danger versus what lowers external danger. For instance, we might consider that developing massive amounts of weaponry can prove useful to stopping certain kinds of external Departure Points, (such as asteroids or whatnot), however this becomes a net negative when we consider the prospect of a biosphere turning that weaponry back on itself in the form of war. In effect, by taking one measure to diminish the external danger, we have actually increased the internal danger. This is one of the first problems that a developing Golden Biosphere will have to contend with. And it is one that it would have solved at some point in its history through essentially an optimization procedure. However, of course optimization is not inherently sufficient here unless its fixed point tends to zero -- and as such perhaps we can consider a Golden Biosphere as A) optimizing its internal and external departure precautions against one another, while B) progressively lowering the sum of all dangers in the process. This will be elaborated on further in the Technoanatomy section, but first we need to explicitly define Aurutance and Golden Scores. 

\subsection{Golden Scores}

Considering how the fundamental ideal of the future is to operate by the Golden Initiative, it follows that we can measure such an adherence. Specifically, given that the Golden Answer directly seeks to have our biosphere live forever, we might consider for this the probability of such a thing occurring. Consider some decision that a person has to make, either adding state $A$ to their Aethus, or instead adding the mutually exclusive state $B$. It follows from the Golden Initiative that they ought to choose the answer with the highest probability of making the biosphere survive forever relative to their Aethus. However, note that it matters how difficult $A$ might be to pick up rather than $B$. Technically speaking, $A$ is a better choice if and only if the following holds, (where $G$ is the state of the biosphere surviving forever). 
\begin{equation}
    P \! \left( A \right) P \! \left( G \mid A \right) \geq P \! \left( B \right) P \! \left( G \mid B \right)
\end{equation}
That is, the probability of actually choosing the option also factors into its consequence. So, for instance, if $A$ and $B$ were in your Aethus in a 20-80 disagreeing superposition, it would follow that $B$ ought to have an additional factor of four to be multiplied onto it before comparing the probabilities of eternal survival relative to the respective child Aethus of yours. Written out, that particular notion is just an algebraic step on the previous one. 
\begin{equation}
    \frac{P \! \left( G \mid A \right)}{P \! \left( G \mid B \right)} \geq \frac{P \! \left( B \right)}{P \! \left( A \right)}
\end{equation}
Note that due to the infinitesimal state of things at infinity, it would in fact be mathematically better to write this as follows.
\begin{equation}
    \lim_{t \to \infty} \frac{P \! \left( S_t \mid A \right)}{P \! \left( S_t \mid B \right)} \geq \frac{P \! \left( B \right)}{P \! \left( A \right)}
\end{equation}
This is so where $S_t$ simply measures the state of the biosphere surviving until time $t$ in the future. The limit proceeds to take the needed ratio in finite pieces, and then work them out to infinity. 

As for an actual means of measuring such a thing, the best option is probably to make it linear, and that means using a logarithm. Consider the following quantity, being called ``Aurutance." (A portmanteau of the Latin word for gold, and the English word importance).
\begin{equation}
    \lim_{t \to \infty} \ln \! \left( \frac{P \! \left( S_t \mid A \right) P \! \left( A \right)}{P \! \left( S_t \mid B \right) P \! \left( B \right)} \right)
\end{equation}
Following from the properties of logarithms, we see that positive Aurutance can be considered to represent a better decision, whereas negative Aurutance might represent a worse decision. 

Consider, for a moment, the societal perception of Aurutance. Whether or not this comparison will age well is hard to tell at present, but perhaps we might suppose that we draw an analogy between Aurutance in the future, and ``money" now. It is a quantity that is abstract, and yet we assign fundamental value to, often even regarding things like quality of life or importance to the world as being defined through it. With the prospect of whether or not it will be relevant being debatable, perhaps we might imagine treating Aurutance as its own kind of currency of sorts, and then undergoing reasoning with this central lens being in place. Perhaps there is a standardized ``dollar" of sorts to Aurutance, which we can refer to as the ``Golden Score," which, for the sake of running a thought experiment, might be of vital importance to, in essence, ``running" the future. Just as one billion dollars in the 21\textsuperscript{st} century can be used to fund some societal achievement, it follows that we might visualize the very existence of the future itself as being powered by the Golden Score. Every Golden Score taken off the table is a scenario of utmost devastation as life on Earth perishes from somewhere in the disagreeing superposition. As such, the entire mission statement of the Golden Initiative can be thought of as the drive to maximize our Golden Scores, and to do so proficiently. This is something that we can consider to be very prevalent amongst any future descendants of ours. 

And, inductively, it can be seen that we’re already bartering with Golden Scores, even if we don’t yet realize it as a society. As will be further elaborated upon in the Technoanatomy section, people like Elon Musk are generating essentially no Golden Scores for us -- especially compared to what they could be doing if they put their (physical) money in better places. This is comically ironic, because as people like this go around bragging about their net worth, they fail to see how, by Golden reasoning, they are among the poorest people in society\footnote{At least by measure of Golden Scores.}. (Note that this is not applicable to all billionaires -- just primarily narcissists). It could even be argued that Elon Musk and others are putting us into debt with Golden Scores, but this is somewhat more subjective. What is clear, at least, is that they are bad investments when it comes to Golden Scores, and if anything they are deeply inefficient at the production of them. In addition to pointing out their wrongdoings, the Technoanatomy section will discuss what we need to supplement this with in order to make a large profit in Golden Scores. 

\subsection{Ontology of Humans}

A major note that Golden reasoning makes about the ontology of humans, is that it disagrees with the notions of misanthropic antinatalism. 

\begin{prince}[The Golden Security Principle]
    If humans were to wipe ourselves out or otherwise cease our modern activities on this planet out of service to the remainder of the biosphere, that would actually be ultimately doing it a disservice, because said actions would be triggering a Departure Point. (I.e. we would cut off the biosphere’s access from the Golden Point, in essence guaranteeing its eventual death). This is simply due to logic from the Optimal Earth Conjecture \cite{Nexus} that trying to make humans evolve from scratch again, so as to develop the Golden Point, would be an Imperfection Magnitude Drop. 
\end{prince}
This is technically a form of techno-optimism, however in a substantially different light to how techno-optimism is usually displayed. To begin with, Golden reasoning generally not only implies that humans ought to continue to develop the Golden Point, but furthermore that such a thing is our legitimate ontological purpose on Earth, itself following from the Golden Extrapolation. Therefore, through this we ought to consider ourselves as fundamentally agents who serve the biosphere, or most particularly serve the Golden Point. At least that is our direct ontological ``purpose" in this system of life. Again, it cannot be emphasized enough in Golden reasoning how we are not insignificant and we are not independently or aimlessly operating agents. All of these notions are assumed in Postmodernism, and as such it serves as a clear guideline for how our perception of our own ontology should change in the coming decades. 

Regarding classical techno-optimism, 